% !TeX encoding = UTF-8
% Français de France. Charger un seul archétype depuis le document principal.
\newcommand{\cvsummary}{Ingénieur en mécanique et mécatronique formé à NYU Abu Dhabi, avec une expérience en R\&D en sciences spatiales aux ÉAU. Association du traitement de données assisté par intelligence artificielle et apprentissage automatique à une pratique de la commande embarquée, de la conception de capteurs et des systèmes robotiques.}

% --- COMPÉTENCES ---
\newcommand{\skills}{
    \section*{COMPÉTENCES TECHNIQUES}
    {\small
    \textbf{Systèmes autonomes et robotique :} Fusion de capteurs (IMU), commande en temps réel, acquisition de nuages de points 3D et téléopération, traitement d'images multispectrales, MATLAB/Simulink, simulations de vol.

    \vspace{1pt}
    \textbf{Programmation et IA :} Python, C++, MATLAB, JavaScript, PostgreSQL, Git ; NumPy, SciPy, notions de PyTorch/fastai, traitement d'images par apprentissage automatique, bases de vision par ordinateur, traitement automatique du langage et intégration d'API, chaînes de traitement de données.

    \vspace{1pt}
    \textbf{Systèmes embarqués et électronique :} Microcontrôleurs, PWM, ADC/DAC, codeurs en quadrature, commande de moteurs à courant continu, machines à états finis, Arduino, Raspberry Pi, automates programmables, IoT, intégration matériel-logiciel.

    \vspace{1pt}
    \textbf{Structures aérospatiales et simulation :} CATIA V5, SolidWorks (CSWA), ANSYS (Mechanical, Fluent), COMSOL Multiphysics, analyse par éléments finis (FEA), mécanique des fluides numérique (CFD), analyse vibratoire, conception allégée, matériaux composites, fabrication additive.

    \vspace{1pt}
    \textbf{Microfabrication et MEMS :} Conception et fabrication de capteurs MEMS piézorésistifs, photolithographie et dessin de masques, dépôt de couches minces (PVD/pulvérisation cathodique), gammes de fabrication à l'échelle du wafer, sélection de matériaux, métrologie et cotation dimensionnelle et géométrique (GD\&T).

    \vspace{1pt}
    \textbf{Données spatiales et géospatiales :} Cartographie planétaire, données satellitaires et multispectrales, PostGIS, numérisation 3D.

    \vspace{1pt}
    \textbf{Langues :} Espagnol langue maternelle, anglais C2, français C2, italien A1, allemand A1.
    }
}

% --- PROJETS ---
\newcommand{\vawttitle}{VAWT -- Récupération d'énergie éolienne urbaine}
\newcommand{\vawtdates}{En cours (date de début non précisée)}
\newcommand{\vawtlocation}{Abou Dabi, ÉAU}
\newcommand{\vawtdescOne}{Développement d'un concept de récupération d'énergie éolienne urbaine par conception d'une éolienne à axe vertical ciblant les écoulements variables générés par le passage des véhicules.}
\newcommand{\vawtdescTwo}{Avancement de l'évaluation aérodynamique par simulations CFD sous ANSYS Fluent, en modélisant le concept d'éolienne avant les décisions de conception physique.}

\newcommand{\fsaetitle}{Formula SAE -- Suspension et châssis en fibre de carbone}
\newcommand{\fsaedates}{Automne 2024 -- Aujourd'hui}
\newcommand{\fsaelocation}{Abou Dabi, ÉAU}
\newcommand{\fsaedescOne}{Développement et réalisation d'un ensemble structurel de véhicule comprenant suspension et châssis en fibre de carbone, par modélisation CAO et analyse par éléments finis. Participation à Formula SAE à l'Universidad de Oviedo et à NYUAD.}
\newcommand{\fsaedescTwo}{Évaluation du comportement de la suspension par analyses cinématiques sous SolidWorks et ANSYS afin d'étudier les performances dynamiques.}

\newcommand{\missiletitle}{Simulateur de dynamique d'interception de missiles}
\newcommand{\missilelocation}{Besançon, France}
\newcommand{\missiledescOne}{Développement sous MATLAB d'une simulation 3D d'interception de missiles avec loi de guidage par navigation proportionnelle, intégration numérique et visualisation des trajectoires en temps réel avec détection d'impact.}

\newcommand{\tactitle}{Mini-réfrigérateur -- Système thermoélectrique TEC à deux étages}
\newcommand{\tacdates}{Printemps 2025}
\newcommand{\taclocation}{Abou Dabi, ÉAU}
\newcommand{\tacdescOne}{Développement d'un dispositif thermique de refroidissement thermoélectrique à deux étages sans fluide frigorigène, grâce à des modules TEC à l'état solide remplaçant le circuit frigorifique.}
\newcommand{\tacdescTwo}{Intégration de capteurs et d'une instrumentation thermique en temps réel dans le prototype pour suivre le comportement du système de refroidissement.}

\newcommand{\orbitaltitle}{OrbitalGrid -- Système SIG pour applications spatiales}
\newcommand{\orbitaldates}{06/2024 - Aujourd'hui}
\newcommand{\orbitallocation}{Abou Dabi, ÉAU}
\newcommand{\orbitaldescOne}{Développement d'un système de données géospatiales à travers un projet de cartographie aérospatiale fondé sur SQL, intégrant numérisation 3D et données satellitaires.}
\newcommand{\orbitaldescTwo}{Mise en place d'une gestion de données spatiales sous PostgreSQL, en structurant les informations géospatiales pour des applications cartographiques.}

\newcommand{\softrobtitle}{Exosquelette robotique pour la rééducation après un AVC}
\newcommand{\softrobdates}{04/2024 - 05/2024}
\newcommand{\softroblocation}{Abou Dabi, ÉAU}
\newcommand{\softrobdescOne}{Conception d'un mécanisme de robotique souple pour la rééducation, sous la forme d'un exosquelette de doigts inspiré de l'origami et actionné par fils en alliage à mémoire de forme.}
\newcommand{\softrobdescTwo}{Intégration de matériaux souples et d'éléments actifs dans un assemblage de polymère Dragon Skin et de fils SMA, conçu pour reproduire la fonction musculaire des doigts.}

\newcommand{\exoplanettitle}{Simulateur de transits d'exoplanètes}
\newcommand{\exoplanetdates}{Printemps 2025}
\newcommand{\exoplanetlocation}{Abou Dabi, ÉAU}
\newcommand{\exoplanetdescOne}{Développement en Python d'un simulateur physique de transits d'exoplanètes, par implémentation des relations de mécanique orbitale et de physique stellaire.}
\newcommand{\exoplanetdescTwo}{Génération de données d'observation synthétiques sous forme de séries temporelles de courbes de lumière, par calcul des transits à partir de paramètres astrophysiques.}

\newcommand{\windTurbinetitle}{Éolienne flottante en mer -- Stabilisation active}
\newcommand{\windTurbinedates}{05/2026 - 06/2026}
\newcommand{\windTurbinelocation}{Besançon, France}
\newcommand{\windTurbinedescOne}{Stabilisation en boucle fermée de la nacelle en laboratoire lors d'essais en bassin à houle, grâce à une commande PID à deux axes sous Arduino avec retour d'orientation d'un MPU6050.}
\newcommand{\windTurbinedescTwo}{Conception mécanique d'une plateforme offshore semi-submersible à trois colonnes avec un système de commande pour stabiliser un prototype d'éolienne à l'échelle 1:1000, en analysant le tangage et l'inclinaison au moyen d'un capteur MPU6050.}

\newcommand{\carbonmetertitle}{CARBONMETRE -- Écoconception et mesure carbone}
\newcommand{\carbonmeterdates}{09/2025 - 12/2025}
\newcommand{\carbonmeterlocation}{Besançon, France}
\newcommand{\carbonmeterdescOne}{Application de l'ingénierie de la mesure à l'évaluation environnementale dans un projet de mesure carbone combinant métrologie, fabrication et écoconception.}
\newcommand{\carbonmeterdescTwo}{Mise en relation de la fabrication et de l'analyse de durabilité par application des principes d'ingénierie de fabrication industrielle au projet.}

\newcommand{\hvactitle}{Amélioration thermique d'une chambre en résidence -- Refroidissement passif et évaporatif}
\newcommand{\hvacdates}{06/2026}
\newcommand{\hvaclocation}{Besançon, France}
\newcommand{\hvacdescOne}{Réduction déclarée de température jusqu'à 3 \textdegree{}C par combinaison d'un écran solaire extérieur et d'un prototype de refroidissement à eau gelée et ventilateur dans une chambre de 9 m$^2$.}
\newcommand{\hvacdescTwo}{Réduction de l'exposition solaire directe aux heures de pointe de 14 h à 21 h, par superposition à l'extérieur de papier aluminium réfléchissant, de PSE de 1 cm et de carton ondulé.}

\newcommand{\controlBoxtitle}{Installation électrique complète -- Équipements pour événements de taille moyenne}
\newcommand{\controlBoxdates}{06/2023 - 08/2023}
\newcommand{\controlBoxlocation}{Mexico, Mexique}
\newcommand{\controlBoxdescOne}{Réduction de 30\% du temps de montage du système par conception d'un coffret de commande de 12 entrées pour des installations électriques événementielles de taille moyenne.}
\newcommand{\controlBoxdescTwo}{Développement d'un système électrique coordonné à automate programmable pour l'événementiel, en intégrant les fonctions de commande d'un environnement d'exploitation en direct.}

\newcommand{\hyperlooptitle}{NYC Hyperloop -- Véhicule à échelle réduite}
\newcommand{\hyperloopdates}{01/2024 - 05/2024}
\newcommand{\hyperlooplocation}{New York, États-Unis}
\newcommand{\hyperloopdescOne}{Pilotage du développement mécanique d'un système de transport avancé en tant que responsable du sous-groupe mécanique, par codéveloppement d'un Hyperloop à échelle réduite pour un concept de liaison New York--Long Island.}
\newcommand{\hyperloopdescTwo}{Évaluation du transport électromagnétique par simulations de lévitation magnétique et utilisation de COMSOL, ANSYS et SolidWorks pour la conception du système.}

\newcommand{\dronetitle}{Quadricoptère fonctionnel -- Structure allégée}
\newcommand{\dronedates}{03/2024 - 05/2024}
\newcommand{\dronelocation}{New York, États-Unis}
\newcommand{\dronedescOne}{Réduction de 40\% du poids du quadricoptère par conception d'une structure imprimée en 3D en PLA/TPU et intégration dans un véhicule fonctionnel.}
\newcommand{\dronedescTwo}{Évaluation de la résistance aux chocs de la structure par essais de chute comparés aux critères de performance pendant le développement.}

\newcommand{\trentxwbtitle}{Rolls-Royce Trent XWB -- Modèle 3D fonctionnel}
\newcommand{\trentxwbdates}{01/2024 - 05/2024}
\newcommand{\trentxwblocation}{New York, États-Unis}
\newcommand{\trentxwbdescOne}{Codéveloppement d'un modèle détaillé du système de propulsion Trent XWB, avec représentation 3D fonctionnelle complète sous SolidWorks et Fusion 360 et évaluation structurelle et des écoulements sous ANSYS.}
\newcommand{\trentxwbdescTwo}{Évaluation du comportement aérodynamique par simulations et outils d'analyse des écoulements sous ANSYS.}

\newcommand{\liostitle}{LIOS -- Assistant vocal avec IA}
\newcommand{\liosdates}{12/2024 - Aujourd'hui}
\newcommand{\lioslocation}{Abou Dabi, ÉAU}
\newcommand{\liosdescOne}{Développement d'un système interactif avec IA sous la forme d'un assistant vocal Python associant traitement automatique du langage, données en temps réel et intégration d'API.}
\newcommand{\liosdescTwo}{Mise en place d'interactions tenant compte du contexte grâce à une mémoire contextuelle conservant les informations pertinentes entre les tâches de l'assistant.}

\newcommand{\fitpaltitle}{FitPal -- Dispositif intelligent d'entraînement physique}
\newcommand{\fitpaldates}{09/2021 - 12/2021}
\newcommand{\fitpallocation}{Abou Dabi, ÉAU}
\newcommand{\fitpaldescOne}{Codéveloppement d'un produit de mesure embarqué fondé sur M5Stack : un dispositif intelligent d'entraînement physique pour l'analyse d'exercices en temps réel.}
\newcommand{\fitpaldescTwo}{Connexion du matériel à une interface mobile par fonctions de transfert de données IoT et intégration d'Arduino IoT App.}

\newcommand{\covidtitle}{Cabine de désinfection automatisée}
\newcommand{\coviddates}{09/2020 - 02/2021}
\newcommand{\covidlocation}{Mexico, Mexique}
\newcommand{\coviddescOne}{Développement d'un prototype automatisé de cabine de désinfection sous Arduino, intégrant électronique de commande et système de brumisation.}
\newcommand{\coviddescTwo}{Mise en place d'une communication distante par ajout de liaisons sans fil au système de commande du prototype.}

\newcommand{\rovertitle}{Rover téléopéré -- Numérisation spatiale 3D}
\newcommand{\roverdates}{09/2026 - Aujourd'hui}
\newcommand{\roverlocation}{Oviedo, Espagne}
\newcommand{\roverdescOne}{Développement en cours d'une plateforme mobile de mesure combinant commande embarquée, acquisition de nuages de points 3D et visualisation à distance, avec la numérisation spatiale téléopérée de bout en bout comme objectif d'intégration.}
\newcommand{\roverdescTwo}{Développement en cours de la téléopération du rover par intégration de la commande embarquée à la plateforme mobile, avec le fonctionnement télécommandé comme jalon cible.}

\newcommand{\microsensortitle}{Microcapteur piézorésistif de pression de détonation}
\newcommand{\microsensordates}{05/2026 - 06/2026}
\newcommand{\microsensorlocation}{Besançon, France}
\newcommand{\microsensordescOne}{Réduction des itérations de dimensionnement du capteur par automatisation des calculs géométriques et de masques dans un outil Python imposant directement des contraintes de K, $\Delta R/R_0$ et résistivité.}
\newcommand{\microsensordescTwo}{Pilotage de la sélection de matériaux sous contraintes pour un capteur à l'échelle du wafer répondant à une interface de 50 $\Omega$ et à une tenue de 10 GPa, en conciliant exigences électriques, mécaniques et économiques.}

% --- EXPÉRIENCE ---
\newcommand{\carbonorg}{Energy Innovation Capital}
\newcommand{\carbontitle}{Analyste de marché -- Captage du carbone}
\newcommand{\carbondates}{10/2023 - 01/2024}
\newcommand{\carbonlocation}{Abou Dabi, ÉAU}
\newcommand{\carbonbullets}{%
  \item Appui aux décisions d'investissement par recommandations sur des jeunes entreprises de captage du carbone, grâce à l'analyse de technologies CCUS pour un fonds de capital-risque d'Abou Dabi.
  \item Transformation d'avancées scientifiques et techniques du captage du carbone en travaux de recherche exploitables et en présentations pour les analystes seniors.
  \item Clarification des incertitudes commerciales par analyse des risques d'investissement, des tendances de marché et des moteurs technologiques du secteur CCUS.
  \item Mise en relation de la recherche en ingénierie et de la stratégie d'investissement par recommandations combinant résultats techniques et analyse de données de marché.
}

\newcommand{\spacelaborg}{NYU Abu Dhabi, Centre des sciences spatiales}
\newcommand{\spacelabtitle}{Ingénieur et chercheur R\&D -- Plateforme de sciences planétaires}
\newcommand{\spacelabdates}{05/2025 - 09/2025}
\newcommand{\spacelablocation}{Abou Dabi, ÉAU}
\newcommand{\spacelabbullets}{%
  \item Amélioration de 40\% de l'efficacité de la chaîne de traitement des données par optimisation de la compression et de la reprojection, avec correction d'images assistée par apprentissage automatique sur des jeux de données à l'échelle planétaire.
  \item Déploiement d'une plateforme de cartographie scientifique sous forme d'une \href{https://romvio.github.io/mars-viewer/}{application web et mobile}, par intégration d'images multispectrales d'Emirates Mars Mission HOPE avec CSS/JavaScript ; soumission ultérieure à l'App Store d'Apple aux ÉAU, étendant l'atlas planétaire au-delà de sa version web.
  \item Intégration d'observations multi-images dans des cartes planétaires par calibration des recouvrements et implémentation des projections cylindrique et Robinson.
}

\newcommand{\geoorg}{Global Education NYUAD}
\newcommand{\geotitle}{Chargé de coordination Global Education -- Étudiants en mobilité entrante}
\newcommand{\geodates}{12/2025 - 05/2025}
\newcommand{\geolocation}{Abou Dabi, ÉAU}
\newcommand{\geobullets}{%
  \item Organisation de programmes interculturels aux ÉAU, dont des sorties universitaires et culturelles pour étudiants en mobilité entrante, par coordination d'activités internationales à NYU Abu Dhabi.
  \item Coordination des opérations par encadrement de l'équipe de coordinateurs et planification logistique pour des groupes étudiants variés.
  \item Contribution à la sécurité des groupes par planification des protocoles et intégration de précautions opérationnelles à l'organisation des sorties.
  \item Maintien de la communication entre parties prenantes en tant que référent des étudiants en mobilité, en les mettant en relation avec les services d'accompagnement universitaires.
}

\newcommand{\geoOneorg}{Global Education NYUAD}
\newcommand{\geoOnetitle}{Chargé de coordination Global Education -- Étudiants de première année}
\newcommand{\geoOnedates}{09/2026 - 12/2024}
\newcommand{\geoOnelocation}{Abou Dabi, ÉAU}
\newcommand{\geoOnebullets}{%
  \item Accompagnement de l'intégration en première année par séances d'accueil et présentations structurées de la vie universitaire.
  \item Conseil continu entre pairs par mentorat et accompagnement pendant la transition vers l'université.
  \item Transmission d'informations sur les installations et ressources universitaires lors de visites du campus pour les nouveaux étudiants.
  \item Contribution à l'intégration communautaire par planification d'événements universitaires et culturels inclusifs pour les étudiants de première année.
}

\newcommand{\hcorg}{Centre de santé de NYUAD}
\newcommand{\hctitle}{Emploi étudiant -- Responsable des réseaux sociaux}
\newcommand{\hcdates}{09/2024 - 05/2025}
\newcommand{\hclocation}{Abou Dabi, ÉAU}
\newcommand{\hcbullets}{%
  \item Promotion des ressources de bien-être par campagnes de santé étudiantes ciblées et adaptation des messages au public de NYU Abu Dhabi.
  \item Production de supports de communication numérique et de contenus publiés sur les réseaux sociaux pour promouvoir les initiatives de santé.
  \item Rôle de point de contact des étudiants, avec prise en charge des demandes entrantes et coordination des réponses du centre de santé.
  \item Contribution à la résolution des demandes en associant communication et accompagnement direct des étudiants.
}

\newcommand{\alpesorg}{Alpes Technologies México}
\newcommand{\alpestitle}{Ingénieur projets junior -- Gestion d'entrepôt et systèmes électriques}
\newcommand{\alpesdates}{06/2024 - 09/2024}
\newcommand{\alpeslocation}{Mexico, Mexique}
\newcommand{\alpesbullets}{%
  \item Amélioration de 60\% de l'efficacité opérationnelle et de fabrication en remplaçant la gestion manuelle des stocks par un système de gestion d'entrepôt sur mesure en HTML/JavaScript et Google Drive.
  \item Mise en place d'une traçabilité industrielle par suivi des numéros de lot et mises à jour bidirectionnelles des stocks en temps réel, avec lecture QR, accès par rôle et rapports automatisés.
  \item Amélioration des performances électriques industrielles, avec facteur de puissance jusqu'à 99\%, réduction de 15\% des coûts énergétiques et de 25\% des arrêts, par analyse des postes et installation de batteries de condensateurs.
  \item Diagnostic électrique après l'incendie d'un complexe hôtelier par mesures sur site, avec schémas unifilaires et recommandations correctives présentés à la direction générale.
}

\newcommand{\epigsaorg}{EPIGSA}
\newcommand{\epigsatitle}{Assistant en ingénierie des données -- SIG et hydraulique}
\newcommand{\epigsadates}{05/2024 - 07/2024}
\newcommand{\epigsalocation}{Mexico, Mexique}
\newcommand{\epigsabullets}{%
  \item Réduction de 25\% des erreurs de données par automatisation du contrôle qualité des données d'infrastructures municipales de traitement des eaux.
  \item Intégration des informations géospatiales d'infrastructures par consolidation de données SIG validées sous PostgreSQL pour les projets de traitement des eaux.
  \item Amélioration de la fiabilité des données d'ingénierie par validation SIG et vérification des informations d'infrastructure lors de leur intégration en base.
  \item Mise en relation des bases de données et des infrastructures physiques par intégration spatiale et contrôle qualité automatisé pour des projets de traitement des eaux sous PostgreSQL.
}

\newcommand{\solarshiporg}{Solar Ship}
\newcommand{\solarshiptitle}{Stagiaire en stratégie et conception de mobilité aérienne}
\newcommand{\solarshipdates}{05/2023 - 08/2023}
\newcommand{\solarshiplocation}{Toronto, Canada}
\newcommand{\solarshipbullets}{%
  %\item Développement d'un concept d'infrastructure de logistique aérienne pour les zones isolées de Zambie et de Tanzanie, en répondant aux besoins d'accès de communautés hors réseau.
  \item Contribution au développement de systèmes aérospatiaux électriques dans l'initiative Peace+Freedom, par conception d'un système aéronautique entièrement électrique pour la connectivité de communautés isolées de Zambie et de Tanzanie hors réseau.
  \item Mise en relation de la conception du véhicule et du contexte d'exploitation dans une application de logistique isolée, en considérant les infrastructures avec le système aérien.
  \item Intégration d'exigences techniques et sociales dans un projet multinational de mobilité durable, en associant conception aérospatiale et objectifs d'accès des communautés.
}

\newcommand{\crousorg}{CROUS Besançon Bourgogne Franche-Comté}
\newcommand{\croustitle}{Emploi étudiant -- Restauration universitaire}
\newcommand{\crousdates}{10/2024 - Aujourd'hui}
\newcommand{\crouslocation}{Besançon, France}
\newcommand{\crousbullets}{%
  \item Maintien de la continuité du travail en cuisine aux heures de pointe par contribution à un service de restauration universitaire à fort volume.
  \item Respect de la discipline des procédures par nettoyage, hygiène de cuisine et tâches de sécurité alimentaire.
  \item Contribution à la continuité opérationnelle par réapprovisionnement et assistance pendant le service de restauration.
  \item Participation coordonnée aux tâches de cuisine et de service, avec adaptation aux besoins opérationnels soumis à des délais courts.
}

\newcommand{\rallyorg}{Championnat du monde des rallyes -- Guanajuato}
\newcommand{\rallytitle}{Commissaire technique officiel FIA}
\newcommand{\rallydates}{11/2019 - 04/2020}
\newcommand{\rallylocation}{Guanajuato, Mexique}
\newcommand{\rallybullets}{%
  \item Réalisation de contrôles techniques réglementaires et du marquage officiel des pneumatiques au Championnat du monde des rallyes FIA, selon les procédures de l'épreuve.
  \item Contribution à la conformité des véhicules de compétition par vérifications techniques au regard des exigences réglementaires.
  \item Coordination des inspections avec les commissaires responsables et les pilotes dans les délais critiques d'une compétition sous forte pression.
  \item Application de procédures de marquage et de vérification dans des travaux de contrôle décrits conformément aux normes ISO 9001.
}

\newcommand{\papaloteorg}{Papalote Museo del Niño}
\newcommand{\papalotetitle}{Chargé d'accueil et de médiation}
\newcommand{\papalotedates}{05/2019 - 07/2019}
\newcommand{\papalotelocation}{Mexico, Mexique}
\newcommand{\papalotebullets}{%
  \item Traduction de concepts techniques pour des non-spécialistes lors de visites d'expositions scientifiques et technologiques, avec explications adaptées aux enfants et aux adultes.
  \item Proposition d'améliorations concrètes de l'expérience de visite, mises en œuvre pendant le travail au musée.
  \item Accompagnement de visiteurs de différents âges en adaptant la communication à leurs besoins.
  \item Mise en relation des contenus techniques et de la participation du public par médiation pratique et explications scientifiques accessibles dans les expositions.
}